\documentclass{article}
\usepackage{amsthm}
\usepackage{comment}
\usepackage{amsmath}
\usepackage{amsfonts}
\usepackage{sectsty}
\usepackage{hyperref}
\usepackage{fancyhdr}
\usepackage[top=1in,bottom=1in,left=1in,right=1in]{geometry}
\usepackage{float}
\usepackage{graphicx}
\usepackage{tabularx}
\usepackage{mathtools}

\newtheorem{theorem}{Theorem}
\newtheorem{axiom}{Axiom}

\title{Lecture 2 - Induction and Cardinality}
\author{Matthew Farah, \href{mailto:farahm11@mcmaster.ca}{$\langle$farahm11@mcmaster.ca$\rangle$}, 400468251}
\date{\today}

\hypersetup{
        colorlinks=true,
        linkcolor=blue,
        filecolor=magenta,
        urlcolor=blue,
        citecolor=blue,
	anchorcolor=blue
}

\fancyhead{}
\fancyhead[L]{Matthew Farah, \href{mailto:farahm11@mcmaster.ca}{$\langle$farahm11@mcmaster.ca$\rangle$}, 400468251}
\fancyhead[R]{Lecture 2 - Induction and Cardinality}

\newenvironment{directsubsubs}{
  \makeatletter
  \renewcommand\thesubsubsection{\thesection.\Roman{subsubsection}}
  \makeatother
}{
  \makeatletter
  \renewcommand\thesubsubsection{\thesubsection.\Roman{subsubsection}}
  \makeatother
}

\renewcommand{\thesubsection}{\thesection.\alph{subsection}}
\renewcommand{\thesubsubsection}{\thesection.\alph{subsection}.\Roman{subsubsection}}
\makeatletter

\sectionfont{\fontsize{12}{12}\bfseries}
\subsectionfont{\fontsize{10}{12}\normalfont}
\subsubsectionfont{\fontsize{10}{12}\normalfont}

\newcolumntype{Y}{>{\centering\arraybackslash}X}
\renewcommand{\arraystretch}{1.8}

\newcommand{\m}[1]{\ensuremath{\displaystyle #1}}

\begin{document}

\maketitle
\pagestyle{fancy}

\begin{theorem}
The set $\wp(\mathbb{Q})$ is uncountable.
\end{theorem}

\begin{proof}
	We begin by restating the corollary found in the previous lecture that $\mathbb{Q}$ is denumerable. Since $\mathbb{Q}$ is denumerable, we can construct an ordered sequence of all the elements of $\mathbb{Q}$ in the form $(r_1, r_2, \ldots, r_i, \ldots )$. Let $S = \{s (s_i)\: s_i \in \{0, 1\}\}$ represent the set of all sequences of zeros and ones. Letting $E \subseteq Q$, consider the function:
	
	\begin{align*}
		f(E) = \begin{cases} 1, \; \text{if} \; r_i \in E \\ 0, \; \text{if} \; r_i \notin E \end{cases}
	\end{align*}

	We can see that this function is injective, because every sequence in $S$ has some corresponding subset in $\mathbb{Q}$ and surjective, since two different subsets of $E_1, E_2 \subseteq \mathbb{Q}, \; E_1 \neq E_2$ must have at least a single element not in common between them, and so necessarily produce different sequences. Therefore, $f$ is a bijective function with a domain of all the subsets of $\mathbb{Q}$ $(\wp(\mathbb{Q}))$ and a co-domain of all possible sequences of zeros and ones. Thus, we know that $\wp(\mathbb{Q})$ and $S$ must have the same size. Now we know that if we can show that $S$ is uncountable, then we have also successfully shown $\wp(\mathbb{Q})$ is uncountable.

	To show $S$ is uncountable, we proceed with an argument from contradiction. Assume, for the purposes of contradiction that $S$ is indeed denumerable. Then, just as we did with $\mathbb{Q}$, we can construct an ordered sequence of all elements of $S$. But, then we can represent $S$ with the following table.
	
	\vspace{1cm}

	\begin{tabularx}{\linewidth}{ c || Y  Y  Y  Y  Y  Y  Y |}
		$s \in S$ & $s_1$ & $s_2$ & $s_3$ & $\ldots$ & $s_i$ & $\ldots$ \\
		\hline
		$s^{(1)}$ & $s_1^1$ & $s_2^1$ & $s_3^1$ & $\ldots$ & $s_i^1$ & $\ldots$ \\
		$s^{(2)}$ & $s_1^2$ & $s_2^2$ & $s_3^2$ & $\ldots$ & $s_i^2$ & $\ldots$ \\
		$s^{(3)}$ & $s_1^3$ & $s_2^3$ & $s_3^3$ & $\ldots$ & $s_i^3$ & $\ldots$ \\
		$\ldots$ & $\ldots$ & $\ldots$ & $\ldots$ & $\raisebox{-1ex}{$\ddots$}$ & $\ldots$ & $\ldots$ \\
		$s^{(i)}$ & $s_1^i$ & $s_2^i$ & $s_3^i$ & $\ldots$ & $s_i^i$ & $\ldots$ \\
		$\ldots$ & $\ldots$ & $\ldots$ & $\ldots$ & $\ldots$ & $\ldots$ & $\raisebox{-2.5ex}{$\ddots$}$
	\end{tabularx}

	\vspace{1cm}

	Now consider creating a sequence $(s_k)$ along the diagonal of the table, where for each entry on the diagonal, we simply choose the opposite of whatever its value is (for example, if $S_1^1 = 0, s_k^1 = 1$ and vice versa). We can formalize this by defining the sequence as:

	\begin{align*}
		(s_k) = (s_k^i = 1 - s_i^i)
	\end{align*}

	We therefore have a sequence $(s_k)$ which consists of only zeros and ones (meaning it belongs to $S$) but differs from each entry in the table by exactly one value. This means that our table, which was made under the assumption that we were able to represent \underline{ALL} elements in the set $S$ (because we assumed $S$ to be denumerable) does not contain $(s_k)$, which as stated earlier is in $S$.

	Thus, by contradiction, $S$ cannot be denumerable, and since we already determined that $S$ and $\wp(\mathbb{Q})$ are of the same size, $\wp(\mathbb{Q})$ cannot be denumerable and is thus uncountable.
\end{proof}

\vspace{1cm}

\begin{itemize}
	\item An operator, such as $+$, is commutative if for any two elements of a space $a, b$: $a + b = b + a$
	\item An operator, such as $+$, is associative if for any three elements of a space $a, b, c$: $(a + b) + c = a + (b + c)$
	\item A field $\mathbb{F}$ is a set which possesses certain additive, multiplicative, and distributive properties given below:
	\begin{itemize}
		\item Additive properties:
		\begin{enumerate}
			\item An addition operator $(+)$ which is associative and commutative.
			\item An additive identity element $(0)$ such that when added to an element $a \in \mathbb{F}$, $a + 0 = a$.
			\item All elements in $a \in \mathbb{F}$ have an additive inverse $-a \in \mathbb{F}$ such that $a + -a = 0$.
		\end{enumerate}
		\item Multiplicative properties:
		\begin{enumerate}
			\item A multiplication operator $(\cdot)$ which is associative and commutative.
			\item A multiplicative identity $(1)$ such than when added to an element $a \in \mathbb{F}$, $a \cdot 1 = 1$.
			\item All elements in $a \in \mathbb{F}$ have a multiplicative inverse (besides the additive identity) $\frac{1}{a} \in \mathbb{F}$ such that $a \cdot \frac{1}{a} = 1$.
		\end{enumerate}
		\item Distributive properties:
		\begin{enumerate}
			\item For all elements $a, b, c \in \mathbb{F}$, multiplication is distributive over addition $a \cdot (b + c) = a \cdot b + a \cdot c$.
		\end{enumerate}
	\end{itemize}

\end{itemize}

\begin{axiom}
	$\mathbb{R}$ is a field.
\end{axiom}

\begin{itemize}
	\item This can actually be proven, but for the purpose of this course, this is assumed to be true.
\end{itemize}

\end{document}
